Assume knowledge of 403 and 531 basically. Dunford \& Schwartz. Recall the definition of a \textit{normed \hyperref[lin-def:vector-space]{vector space}}. It induces a metric and hence a topology. A metric space is a \hyperref[mv-def:banach-space]{Banach space} if it is complete. It's often that if you have infinite-dim vector spaces, you have a lot of topologies to choose from, each of which have different roles. We will find out how to utilize each of them. Some examples are $\mathbb{R}^n$ with the $p$-norm, $\ell^p$ such that $\big( \sum_{j=1}^\infty |x_j|^p \big) < \infty$, or $L^p (\mathbb{R}^n)$. 

You should be very careful with your intuition in an infinite-dimensional \hyperref[lin-def:vector-space]{vector space}. You should develop this intuition on what is or isn't obvious in an infinite-dimensional vector space, so we're going over some seemingly ``obvious'' results. 

